\documentclass[11pt]{amsart}

\usepackage[T1]{fontenc}
\usepackage{lmodern}
\usepackage{amsmath,amssymb,amsthm}
\usepackage[margin=2.6cm]{geometry}
\usepackage{booktabs}
\usepackage{longtable,array}
\usepackage[hidelinks]{hyperref}
\usepackage{xurl}

\newtheorem{theorem}{Theorem}[section]
\newtheorem{lemma}[theorem]{Lemma}
\newtheorem{proposition}[theorem]{Proposition}
\newtheorem{corollary}[theorem]{Corollary}
\theoremstyle{definition}
\newtheorem{definition}[theorem]{Definition}
\theoremstyle{remark}
\newtheorem{remark}[theorem]{Remark}

\newcommand{\lean}[1]{\textnormal{\nolinkurl{#1}}}
\newcommand{\Lam}{\Lambda}
\newcommand{\F}{\mathbb{F}_2}
\newcommand{\cC}{\mathcal{C}}
\newcommand{\Lev}{L}
\newcommand{\ins}{\operatorname{ins}}
\newcommand{\rot}{\operatorname{rot}}
\DeclareMathOperator{\seam}{Seam}

\title[J03 and J04: finite tensor algebra and the iterated seam]{J03 and J04: finite tensor
algebra of grid complexes\\ and the iterated seam comparison}
\date{2 November 2026}

\begin{document}

\begin{abstract}
We prove the statements of jobs J03 and J04 of the dependency audit: the finite tensor algebra
of native grid complexes (records M005, M609, M610, M615) and the comparison between the iterated
seam grid and the finite tensor of its factors (records M608, M612). All maps are filtered chain
maps with shifts $(0,0)$, and each comparison carries the canonical class of its source to the
canonical class of its target. The proofs are formalized in Lean, and Section~\ref{sec:lean} lists
the corresponding declarations. As in the record of M609 in the audit, the chain map in M609 is
the restriction of one $\Lam$-linear map of the underlying modules, and the constructed map has the
same property. Remark~\ref{rem:linear} explains why this hypothesis is needed for the
construction.
\end{abstract}

\maketitle

\section{Setting}\label{sec:setting}

Throughout, $K$ is a field of characteristic two and $\Lam=K[T,T^{-1}]$. All homology classes
concerning canonical classes are taken with $K=\F$.

\subsection{Native grid complexes}
For a grid diagram $G$ of size $n$ with set of states $S(G)$, the module $C(G)$ is the free
$\Lam$-module with basis $S(G)$, and $\partial_G\colon C(G)\to C(G)$ is the $\Lam$-linear
full rectangle differential in which every $O$ variable is specialized to $T$ (Lean:
\lean{laurentDifferential}). A monomial $T^kx$ with $x\in S(G)$ has Maslov degree $M(x)-2k$ and
filtration coordinates $(-k,\,A(x)-k)$, where $M$ and $A$ are the Maslov and Alexander gradings of
the grid. For $r\in\mathbb{Z}$ let $C_r(G)\subseteq C(G)$ be the $K$-span of the monomials of
Maslov degree $r$, and for $a,b\in\mathbb{R}$ let $\Lev_{a,b}(G)\subseteq C(G)$ be the
$K$-submodule of chains all of whose monomials $T^kx$ satisfy $-k\le a$ and $A(x)-k\le b$. The
\emph{native complex} $\cC(G)$ is the $\mathbb{Z}$-graded complex of $K$-vector spaces
$(C_r(G),\partial_G)$ with the bifiltration $F_{a,b}\cC_r(G)=C_r(G)\cap\Lev_{a,b}(G)$ (Lean:
\lean{laurentMaslovComplex}, \lean{laurentMaslovBifiltration}). A chain map $\varphi$ between two
bifiltered complexes has \emph{shifts} $(s_1,s_2)$ if $\varphi(F_{a,b})\subseteq F_{a+s_1,b+s_2}$
for all $a,b$ (Lean: \lean{Bifiltration.Shifts}).

\subsection{Canonical classes}
For a grid diagram $G$ of size at least two, $\xi_G\in H_0(\cC(G))$ denotes the normalized class of
the project (Lean: \lean{gridCanonicalClass}). For every choice of $O$ variable, it is the unique
nonzero class in the image of the corresponding polynomial homology, and $c_G$ denotes the cycle of $\cC_0(G)$ with nonnegative exponents chosen to
represent it (Lean: \lean{gridCanonicalCycle}). We use the following proved fact.

\begin{proposition}[\lean{gridCanonicalClass_homotopyEquiv}]\label{prop:class-transport}
Let $G$ and $H$ be grid diagrams of sizes at least two, and let $e$ be a homotopy equivalence
from $\cC(G)$ to $\cC(H)$ over $\F$ whose forward map has shifts $(0,s)$ for some real $s$. Then
$e_*(\xi_G)=\xi_H$.
\end{proposition}

\subsection{Finite and binary tensors}
For a finite index type $\iota$ and grid diagrams $(G_i)_{i\in\iota}$, put
$C(G_\bullet)=\bigotimes_{\Lam,i\in\iota}C(G_i)$ with the differential
$\partial=\sum_{i\in\iota}\partial_i$, where $\partial_i$ applies $\partial_{G_i}$ to the $i$th
factor. Since $K$ has characteristic two, no signs occur and $\partial^2=0$. The piece
$C_r(G_\bullet)$ is the $K$-span of the elements $T^j\otimes_ix_i$ with $x_i\in C_{p_i}(G_i)$ and
$\sum_ip_i-2j=r$, and the level $\Lev_{A,B}(G_\bullet)$ is the $K$-span of the elements
$T^j\otimes_ix_i$ with $x_i\in\Lev_{a_i,b_i}(G_i)$, $\sum_ia_i-j\le A$ and $\sum_ib_i-j\le B$. These
define the native finite tensor complex $\cC(G_\bullet)$ and its bifiltration (Lean:
\lean{laurentFiniteMaslovComplex}, \lean{laurentFiniteMaslovBifiltration}). If every $G_i$ has
size at least two, the \emph{canonical tensor class} is
$\xi_{G_\bullet}=[\otimes_ic_{G_i}]\in H_0(\cC(G_\bullet))$ (Lean:
\lean{gridCanonicalFiniteTensorClass}). The binary tensor $\cC(G\otimes H)$ is defined in the same
way from $C(G)\otimes_\Lam C(H)$, and its canonical class is $\xi_{G\otimes H}=[c_G\otimes c_H]$
(Lean: \lean{laurentTensorMaslovComplex}, \lean{gridCanonicalBinaryTensorClass}).

\subsection{Maps induced on pieces}
Each of the three complexes above has chain groups equal to homogeneous $K$-submodules of one
underlying module, a differential induced by one underlying map, and filtration levels induced by
one family of underlying submodules (Lean: \lean{PieceComplexData}).

\begin{lemma}[\lean{PieceComplexData.hom}, \lean{PieceComplexData.hom_shifts}]\label{lem:pieces}
Let $\cC$ and $\cC'$ be complexes of this form with underlying modules $M$ and $M'$, and let
$f\colon M\to M'$ be a $K$-linear map that maps each piece of degree $r$ into the piece of degree
$r$ and satisfies $f\circ\partial=\partial'\circ f$. Then the restrictions of $f$ form a chain map
$f_\#\colon\cC\to\cC'$. If moreover $f(\Lev_{a,b})\subseteq\Lev'_{a,b}$ for all $a,b$, then $f_\#$ has
shifts $(0,0)$. If $f$ is a linear equivalence whose inverse has the same properties, then $f_\#$
is an isomorphism of complexes.
\end{lemma}

\begin{proof}
The restriction $C_{r+1}\to C'_{r+1}$ composed with $\partial'$ equals $\partial$ composed with the
restriction $C_r\to C'_r$ by $f\circ\partial=\partial'\circ f$. The statement on levels is immediate
from the definition of the bifiltrations, and the restrictions of $f^{-1}$ are inverse to those of
$f$.
\end{proof}

\section{Finite tensor algebra (J03)}\label{sec:j03}

\subsection{The singleton tensor (M005, alias M611)}

\begin{proposition}\label{prop:singleton}
Let $\iota$ have exactly one element $i_0$ and let $(G_i)$ be grid diagrams. The map
$c\mapsto\otimes_i c$ induces an isomorphism $\cC(G_{i_0})\cong\cC(G_\bullet)$ with shifts $(0,0)$
in both directions. If $G_{i_0}$ has size at least two, the isomorphism maps $\xi_{G_{i_0}}$ to
$\xi_{G_\bullet}$.
\end{proposition}

\begin{proof}
The map $c\mapsto\otimes_ic$ is the inverse of the $\Lam$-linear equivalence
$\bigotimes_{i\in\iota}C(G_i)\to C(G_{i_0})$ of Mathlib (\lean{PiTensorProduct.subsingletonEquiv}),
which sends $T^j\otimes_ix_i$ to $T^jx_{i_0}$. Since $\sum_ip_i=p_{i_0}$ and
$\sum_ia_i=a_{i_0}$, it maps $C_r(G_\bullet)$ into $C_r(G_{i_0})$ and $\Lev_{A,B}(G_\bullet)$ into
$\Lev_{A,B}(G_{i_0})$, because multiplication by $T^j$ lowers both coordinates by $j$. Its inverse
maps $C_r(G_{i_0})$ and $\Lev_{A,B}(G_{i_0})$ into the corresponding submodules, taking $j=0$. The
equivalence intertwines $\partial$ and $\partial_{G_{i_0}}$ by
\lean{laurentFiniteDifferential_subsingleton}. Lemma~\ref{lem:pieces} gives the isomorphism and the
shifts. The forward map sends the cycle $c_{G_{i_0}}$ to $\otimes_ic_{G_i}$, hence
$\xi_{G_{i_0}}$ to $\xi_{G_\bullet}$.
\end{proof}

\subsection{Reindexing (M610)}

\begin{proposition}\label{prop:reindex}
Let $e\colon\iota\to\iota'$ be a bijection of finite index types and let $(G_i)_{i\in\iota}$ be grid
diagrams. Reindexing induces an isomorphism $\cC(G_\bullet)\cong\cC(G_{e^{-1}(\bullet)})$ with
shifts $(0,0)$ in both directions. If every $G_i$ has size at least two, it maps $\xi_{G_\bullet}$
to $\xi_{G_{e^{-1}(\bullet)}}$.
\end{proposition}

\begin{proof}
Let $R$ be the $\Lam$-linear equivalence of Mathlib (\lean{PiTensorProduct.reindex}), which sends
$T^j\otimes_ix_i$ to $T^j\otimes_{i'}x_{e^{-1}(i')}$. Since
$\sum_{i'}p_{e^{-1}(i')}=\sum_ip_i$, the map $R$ sends generators of $C_r(G_\bullet)$ to generators
of $C_r(G_{e^{-1}(\bullet)})$, and the same holds for levels. Every pure tensor
$\otimes_{i'}y_{i'}$ of the reindexed family equals $R(\otimes_ix_i)$ for the family
$x_{e^{-1}(i')}=y_{i'}$, so $R^{-1}$ maps generators to generators as well. For the differential,
$R(\partial\otimes_ix_i)=\sum_iR(\otimes_k\tilde x^{(i)}_k)$, where $\tilde x^{(i)}$ replaces $x_i$ by
$\partial_{G_i}x_i$. Substituting $i=e^{-1}(i')$ identifies this sum with
$\partial R(\otimes_ix_i)$. Lemma~\ref{lem:pieces} gives the isomorphism and the shifts, and $R$
sends $\otimes_ic_{G_i}$ to $\otimes_{i'}c_{G_{e^{-1}(i')}}$.
\end{proof}

If two families of sized grid diagrams are equal as families of pairs (size, diagram), their
finite tensor complexes are equal, and the identification preserves the bifiltrations and the
canonical tensor classes (Lean: \lean{laurentFiniteMaslovFamilyIso}). This is used for the
transport along size equalities in Section~\ref{sec:j04}.

\subsection{Powers of the cyclic row translation (M615)}

\begin{proposition}\label{prop:rows}
Let $G$ be a grid diagram of size $n$ and $k\ge0$. Relabeling rows by $\sigma^k$, where $\sigma$ is
the cyclic translation of $\{0,\dots,n-1\}$, induces an isomorphism
$\cC(G)\cong\cC(G^{\sigma^k})$ with shifts $(0,0)$ in both directions. If $n\ge2$, it maps $\xi_G$
to $\xi_{G^{\sigma^k}}$, and its inverse maps $\xi_{G^{\sigma^k}}$ to $\xi_G$.
\end{proposition}

\begin{proof}
For $k=1$, the relabeling preserves the full Laurent coefficients, the Maslov degrees and the
Alexander gradings (\lean{laurentCoefficient_relabelRows_finRotate},
\lean{laurentMaslovDegree_relabelRows_finRotate}, \lean{alexander_relabelRows_finRotate}). Since
relabeling by $\sigma^{k+1}$ is relabeling by $\sigma^k$ followed by relabeling by $\sigma$, induction
on $k$ gives the three invariances for $\sigma^k$. The relabeling $C(G)\to C(G^{\sigma^k})$ of states
is $\Lam$-linear, intertwines the differentials and preserves and reflects pieces and levels.
Lemma~\ref{lem:pieces} gives the isomorphism and the shifts in both directions, and
Proposition~\ref{prop:class-transport} gives the statements on classes.
\end{proof}

\subsection{Insertion of one factor and reassociation (M609)}

Let $\kappa$ be a finite index type presented as the image of an injective map
$\iota\to\kappa$ together with one further index $*$ (Lean: \lean{FactorInsertion}). The examples
are $\kappa=\iota\sqcup\{*\}$ realized as $\mathrm{Option}\,\iota$, and $\kappa=\{0,\dots,N\}$ with
$\iota=\{0,\dots,N-1\}$ and $*=N$. For grid diagrams $(G_k)_{k\in\kappa}$, write $G_\iota$ for the
restricted family. The assignment $(\otimes_ix_i,\,y)\mapsto\otimes_kz_k$, where $z_i=x_i$ and
$z_*=y$, is multilinear in the $x_i$ and linear in $y$. By the universal property it defines a
$\Lam$-bilinear map
\[
\ins\colon C(G_\iota)\times C(G_*)\to C(G_\bullet).
\]

\begin{lemma}[\lean{insertLinear_mem_maslovPiece}, \lean{insertLinear_mem_bifiltrationLevel},
\lean{laurentFiniteDifferential_insertLinear}]\label{lem:insert}
For $z\in C_r(G_\iota)$ and $y\in C_q(G_*)$ we have $\ins(z,y)\in C_{r+q}(G_\bullet)$. For
$z\in\Lev_{A,B}(G_\iota)$ and $y\in\Lev_{c,d}(G_*)$ we have
$\ins(z,y)\in\Lev_{A+c,B+d}(G_\bullet)$. For all $z$ and $y$,
\[
\partial\,\ins(z,y)=\ins(\partial z,y)+\ins(z,\partial_{G_*}y).
\]
\end{lemma}

\begin{proof}
On generators, $\ins(T^j\otimes_ix_i,y)=T^j\otimes_kz_k$, and the degrees and coordinates of the
$z_k$ are those of the $x_i$ together with those of $y$. Their sums are therefore
$\sum_ip_i+q$, $\sum_ia_i+c$ and $\sum_ib_i+d$. For the differential, the sum
$\sum_{k\in\kappa}\partial_k$ splits into the terms $k$ in the image of $\iota$, which give
$\ins(\partial z,y)$, and the term $k=*$, which gives $\ins(z,\partial_{G_*}y)$. Both sides are
$\Lam$-linear in $z$, so the identity on pure tensors suffices.
\end{proof}

\begin{definition}[\lean{LaurentGridFiniteMap}]\label{def:laurent-map}
Let $T$ be a grid diagram and $(D_i)_{i\in\iota}$ a family of grid diagrams. A \emph{filtered
$\Lam$-linear map} $f\colon C(T)\to C(D_\bullet)$ is a $\Lam$-linear map that maps $C_r(T)$ into
$C_r(D_\bullet)$, satisfies $f\circ\partial_T=\partial\circ f$, and maps $\Lev_{a,b}(T)$ into
$\Lev_{a,b}(D_\bullet)$ for all $r,a,b$. By Lemma~\ref{lem:pieces} it induces a chain map
$f_\#\colon\cC(T)\to\cC(D_\bullet)$ with shifts $(0,0)$.
\end{definition}

\begin{definition}[\lean{LaurentGridInduced}, \lean{LaurentTensorInduced}]\label{def:induced}
A chain map $\theta\colon\cC(T)\to\cC(D_\bullet)$ is \emph{$\Lam$-induced} if there is a
$\Lam$-linear map $f\colon C(T)\to C(D_\bullet)$ whose restriction to $C_r(T)$ is $\theta_r$ for every
$r$. The same definition applies to chain maps from a binary tensor complex
$\cC(T\otimes H)$, with $f$ defined on $C(T)\otimes_\Lam C(H)$.
\end{definition}

\begin{lemma}[\lean{LaurentGridInduced.exists_laurentGridFiniteMap}]\label{lem:bridge}
Let $\theta\colon\cC(T)\to\cC(D_\bullet)$ be a $\Lam$-induced chain map with shifts $(0,0)$, induced by
$f$. Then $f$ is a filtered $\Lam$-linear map and $f_\#=\theta$.
\end{lemma}

\begin{proof}
For $c\in C_r(T)$, the element $f(c)=\theta_r(c)$ lies in $C_r(D_\bullet)$. For a state $x$, the basis
vector $x$ lies in $C_{M(x)}(T)$, and the chain map identity for $\theta$ in degree $M(x)$ gives
$f(\partial_Tx)=\partial f(x)$. Both $f\circ\partial_T$ and $\partial\circ f$ are $\Lam$-linear and the
states form a $\Lam$-basis of $C(T)$, hence $f\circ\partial_T=\partial\circ f$. A chain
$c\in\Lev_{a,b}(T)$ is the $K$-linear combination of its monomials $T^kx$, and each of these lies in
$\Lev_{a,b}(T)$, because the level is defined by a condition on each monomial. Each monomial is
homogeneous, so the shifts $(0,0)$ of $\theta$ give $f(T^kx)\in\Lev_{a,b}(D_\bullet)$. Summing gives
$f(c)\in\Lev_{a,b}(D_\bullet)$. Finally $f_\#=\theta$, because both restrict $f$.
\end{proof}

\begin{theorem}[M609]\label{thm:reassociation}
Let $T$ and $H$ be grid diagrams of sizes at least two, let $(D_i)_{i\in\iota}$ be grid diagrams of
sizes at least two, and let $D'$ be the family over $\mathrm{Option}\,\iota$ with $D'_*=H$ and
$D'_i=D_i$. Let $\theta\colon\cC(T)\to\cC(D_\bullet)$ be a $\Lam$-induced chain map with shifts
$(0,0)$ and $\theta_*(\xi_T)=\xi_{D_\bullet}$. Then there is a $\Lam$-induced chain map
$\psi\colon\cC(T\otimes H)\to\cC(D'_\bullet)$ with shifts $(0,0)$ and
$\psi_*(\xi_{T\otimes H})=\xi_{D'_\bullet}$.
\end{theorem}

\begin{proof}
By Lemma~\ref{lem:bridge}, $\theta=f_\#$ for a filtered $\Lam$-linear map $f$. More generally, let
$\kappa$ be presented as above and let $f\colon C(T)\to C(G_\iota)$ be a filtered $\Lam$-linear map.
The map $(c,y)\mapsto\ins(f(c),y)$ is $\Lam$-bilinear, so it defines a $\Lam$-linear map
$\Psi_f\colon C(T)\otimes_\Lam C(G_*)\to C(G_\bullet)$ with $\Psi_f(c\otimes y)=\ins(f(c),y)$.
Lemma~\ref{lem:insert} shows that $\Psi_f$ maps the generators $c\otimes y$ of the binary pieces and
levels into the corresponding pieces and levels, using monotonicity of the levels. For the
differential,
\[
\Psi_f(\partial c\otimes y+c\otimes\partial y)=\ins(\partial f(c),y)+\ins(f(c),\partial y)
=\partial\,\Psi_f(c\otimes y).
\]
By Lemma~\ref{lem:pieces}, $\psi=(\Psi_f)_\#$ has shifts $(0,0)$, and $\psi$ is $\Lam$-induced by
$\Psi_f$ (Lean: \lean{LaurentGridFiniteMap.insertTensor},
\lean{LaurentTensorFiniteMap.laurentTensorInduced}).

For the classes, let $y$ be a cycle of degree zero in $\cC(G_*)$. The map $S_y=\ins(-,y)$ is a
chain map $\cC(G_\iota)\to\cC(G_\bullet)$ by Lemma~\ref{lem:insert}, because $\partial y=0$ (Lean:
\lean{insertCycleHom}). For a cycle $c$ of degree zero in $\cC(T)$, the cycles
$\Psi_f(c\otimes y)$ and $S_y(f(c))$ coincide, so
$\psi_*[c\otimes y]=S_{y*}(f_{\#*}[c])$. With $c=c_T$ and $y=c_{G_*}$ this gives
\[
\psi_*(\xi_{T\otimes G_*})=S_{y*}(\xi_{G_\iota})=[\ins(\otimes_ic_{G_i},c_{G_*})]
=[\otimes_kc_{G_k}]=\xi_{G_\bullet}.
\]
The case $\kappa=\mathrm{Option}\,\iota$ and $G=D'$ is the theorem.
\end{proof}

\begin{remark}\label{rem:linear}
The proposal for M609 cited in the job listing of J03 allows an arbitrary chain map
$\theta\colon\cC(T)\to\cC(D_\bullet)$ of complexes of $K$-vector spaces. The formula
$c\otimes y\mapsto\theta(c)\otimes y$ is well defined on the balanced tensor product
$C(T)\otimes_\Lam C(H)$ only if it respects the relation $(Tc)\otimes y=c\otimes(Ty)$, which
requires $\theta(Tc)=T\theta(c)$. The record of M609 in the dependency audit adds the hypothesis
that $\theta$ is $\Lam$-induced and the same conclusion for $\psi$, and Theorem~\ref{thm:reassociation}
is that statement. The final comparison map of the route is a composite
$\theta_{\mathrm{tr}}\circ\varphi$, where $\varphi$ comes from the cobordism and the iterated
comparison $\theta_{\mathrm{tr}}$ of Section~\ref{sec:j04} is applied last. Composition requires no
linearity of $\varphi$, so no other map of the route needs to be $\Lam$-induced. We do not know whether the proposal for an arbitrary chain map holds,
and it is not needed.
\end{remark}

\section{The iterated seam comparison (J04)}\label{sec:j04}

\subsection{The binary seam (M608)}
For grid diagrams $G$ of size $n+1$ and $H$ of size $m+1$ with $X_G(n)=n$ and $X_H(0)=0$, the
shared corner seam grid $\seam(G,H)$ of size $n+m+1$ is \lean{connectedSumSeamGrid}. Assume also
$O_H(m)=0$. The seam contraction gives a $\Lam$-linear projection
$\pi\colon C(\seam(G,H))\to C(G)\otimes_\Lam C(H)$, the corrected projection followed by the inverse
tensor coordinates (Lean: \lean{seamTensorProjection}). It maps levels into levels
(\lean{seamFullContraction_proj_mem_bifiltration_of_last_O},
\lean{seamTensorCoordinatesEquiv_symm_mem_bifiltration}).

\begin{lemma}[M608, \lean{seamParity_of_positivePaths}]\label{lem:m608}
If the seam positive path comparisons hold for the pair of sizes $(n,m)$, then the parity identity
used as the hypothesis of \lean{seamNativeTensorHomotopyEquiv} holds for every such pair $G,H$.
\end{lemma}

\begin{proof}
The comparisons are equivalent to the equality of the retained seam differential with the
tensor coordinate differential over $\F$
(\lean{seamRetainedDifferential_eq_tensorCoordinate_iff_positive_paths}), and this equality is
equivalent to the parity identity (\lean{seamRetainedDifferential_eq_tensorCoordinate_iff_evalOne}).
\end{proof}

Under the parity identity, $\pi$ commutes with the differentials and maps the piece of degree $r$
into the piece of degree $r$, and $\pi_\#$ is the forward map of the homotopy equivalence
\lean{seamNativeTensorHomotopyEquiv}. It satisfies $\pi_{\#*}(\xi_{\seam(G,H)})=\xi_{G\otimes H}$
by \lean{seamNativeTensor_hom_canonicalClass}. The comparisons hold for all sizes by M601
(\lean{seamPositivePathComparisons}), so Lemma~\ref{lem:m608} applies without hypothesis.

\subsection{Prepared families and the iterated seam grid}

\begin{definition}[\lean{SeamPrepared}]\label{def:prepared}
A family $(D_j)_{0\le j<N}$ of grid diagrams of sizes $s_j$ is \emph{prepared} if every $s_j$ is
positive, $X_{D_0}(s_0-1)=s_0-1$ when $N>0$, and $X_{D_j}(0)=0$ and $O_{D_j}(s_j-1)=0$ for every
$j\ne0$.
\end{definition}

For a grid diagram $G$ of size $n+1$, some power $\sigma^k$ of the cyclic row translation satisfies
$\sigma^k(X_G(n))=n$ (\lean{exists_cyclicRows_send} for $n\ge1$, and $k=0$ for $n=0$). We fix one
such power and put $\rot(G)=G^{\sigma^k}$, so that $X_{\rot(G)}(n)=n$ (Lean: \lean{seamLeftGrid}).

Let $(E_j)_{0\le j\le M}$ be grid diagrams of sizes $p_j+1$ with $X_{E_j}(0)=0$ for $j\ge1$. The
\emph{seam chain} is defined by
\[
T_0=E_0,\qquad T_{k+1}=\seam(\rot(T_k),E_{k+1})\quad(0\le k<M)
\]
(Lean: \lean{seamChainGrid}). By induction $T_k$ has size $\sum_{j\le k}p_j+1$. For a prepared family
$(D_j)_{0\le j<N}$ with $N=M+1$, put $p_j=s_j-1$ and let $E_j$ be $D_j$ viewed as a grid of size
$p_j+1$. The \emph{iterated seam grid} of $D$ is $T_M$, transported to the size
$\sum_js_j-(N-1)=\sum_jp_j+1$ (Lean: \lean{iteratedSeamGrid}, \lean{iteratedSeamSize},
\lean{iteratedSeamSize_succ}). For $N=0$ it is the grid of size zero.

\begin{lemma}[\lean{two_le_iteratedSeamSize}]\label{lem:size}
If $N>0$ and $s_j\ge2$ for all $j$, then $\sum_js_j-(N-1)\ge2$.
\end{lemma}

\begin{proof}
We have $\sum_js_j\ge2N$, hence $\sum_js_j-(N-1)\ge N+1\ge2$.
\end{proof}

\begin{lemma}[\lean{exists_seamPrepared_movesTo}]\label{lem:prepare}
Every family $(D_j)$ of grid diagrams of sizes at least two is related factorwise by grid moves to
a prepared family of grid diagrams of sizes at least two.
\end{lemma}

\begin{proof}
For $j=0$, relabel the rows of $D_0$ by a power of $\sigma$ that sends $X_{D_0}(s_0-1)$ to $s_0-1$
(\lean{exists_cyclicRows_send}). Cyclic row relabelings are grid moves
(\lean{cyclicRows_pow_movesTo}). For $j\ne0$, \lean{exists_seamPreparedGrid} gives a grid $D'_j$ of
size $s_j+1$ reached by grid moves with $X_{D'_j}(0)=0$ and $O_{D'_j}(s_j)=0$.
\end{proof}

\subsection{The comparison (M612)}

\begin{theorem}[M612]\label{thm:iterated}
Assume the seam positive path comparisons for all pairs of sizes. Let $N>0$ and let
$(D_j)_{0\le j<N}$ be a prepared family of grid diagrams of sizes at least two, with iterated seam
grid $G_{\mathrm{tr}}$. Then there is a chain map $\theta\colon\cC(G_{\mathrm{tr}})\to\cC(D_\bullet)$
over $\F$ with shifts $(0,0)$ and $\theta_*(\xi_{G_{\mathrm{tr}}})=\xi_{D_\bullet}$. Since the
comparisons are proved (M601), the conclusion holds without hypothesis.
\end{theorem}

\begin{proof}
We first prove, by induction on $M$, the following statement for the seam chain of grid diagrams
$(E_j)_{0\le j\le M}$ of sizes $p_j+1\ge2$ with $X_{E_j}(0)=0$ and $O_{E_j}(p_j)=0$ for $j\ge1$:
there is a filtered $\Lam$-linear map $f_M\colon C(T_M)\to C(E_\bullet)$ with
$f_{M\#*}(\xi_{T_M})=\xi_{E_\bullet}$ (Lean: \lean{exists_seamChain_finiteTensor_map}).

For $M=0$, Proposition~\ref{prop:singleton} gives $f_0$, the inverse singleton equivalence.

For the inductive step, let $f_M\colon C(T_M)\to C(E_{\le M})$ be given for the family
$(E_j)_{j\le M}$, and let $k$ be the chosen power with $\rot(T_M)=T_M^{\sigma^k}$.
\begin{enumerate}
\item The inverse row relabeling $C(\rot(T_M))\to C(T_M)$ is a filtered $\Lam$-linear map of grid
chains (Proposition~\ref{prop:rows}), and its induced chain map is the inverse isomorphism of
Proposition~\ref{prop:rows}. Hence $g=f_M\circ\sigma^{-k}$ is a filtered $\Lam$-linear map
$C(\rot(T_M))\to C(E_{\le M})$, and
\[
g_{\#*}(\xi_{\rot(T_M)})=f_{M\#*}(\xi_{T_M})=\xi_{E_{\le M}}
\]
by Proposition~\ref{prop:rows} and the induction hypothesis.
\item Insertion of the factor $E_{M+1}$ into the last position, with $\kappa=\{0,\dots,M+1\}$ and
$*=M+1$, gives the map $\Psi_g\colon C(\rot(T_M))\otimes_\Lam C(E_{M+1})\to C(E_{\le M+1})$ of the
proof of Theorem~\ref{thm:reassociation}, with
$(\Psi_g)_{\#*}(\xi_{\rot(T_M)\otimes E_{M+1}})=\xi_{E_{\le M+1}}$.
\item The grid $\rot(T_M)$ satisfies $X_{\rot(T_M)}(s)=s$, where $s+1$ is its size, and $E_{M+1}$
satisfies $X_{E_{M+1}}(0)=0$ and $O_{E_{M+1}}(p_{M+1})=0$. The parity identity holds for this pair
by Lemma~\ref{lem:m608}, and $T_{M+1}=\seam(\rot(T_M),E_{M+1})$. Put $f_{M+1}=\Psi_g\circ\pi$. It is
a filtered $\Lam$-linear map because $\pi$ and $\Psi_g$ preserve pieces, levels and differentials.
Its induced chain map is $(f_{M+1})_\#=(\Psi_g)_\#\circ\pi_\#$, as recorded in
\lean{LaurentTensorFiniteMap.precompSeam_hom}. Therefore
\[
f_{M+1\,\#*}(\xi_{T_{M+1}})=(\Psi_g)_{\#*}(\xi_{\rot(T_M)\otimes E_{M+1}})=\xi_{E_{\le M+1}}.
\]
\end{enumerate}

For the theorem, write $N=M+1$ and let $E_j$ be $D_j$ viewed as a grid of size $(s_j-1)+1$. These
grids satisfy the hypotheses above because the family is prepared and $s_j\ge2$. Let $f_M$ be the
map just constructed. Put
\[
\theta=\tau_D\circ f_{M\#}\circ\tau_{\mathrm{tr}}^{-1},
\]
where $\tau_{\mathrm{tr}}\colon\cC(T_M)\to\cC(G_{\mathrm{tr}})$ is the transport along the size
equality and $\tau_D\colon\cC(E_\bullet)\to\cC(D_\bullet)$ is the transport of families. Both
transports are identities after the size equalities are substituted, so they have shifts $(0,0)$,
$\tau_{\mathrm{tr}}^{-1}$ maps $\xi_{G_{\mathrm{tr}}}$ to $\xi_{T_M}$ by
Proposition~\ref{prop:class-transport}, and $\tau_D$ maps $\xi_{E_\bullet}$ to $\xi_{D_\bullet}$.
Shifts add under composition (\lean{Bifiltration.Shifts.comp}), so $\theta$ has shifts $(0,0)$ and
$\theta_*(\xi_{G_{\mathrm{tr}}})=\xi_{D_\bullet}$.
\end{proof}

\begin{remark}
The condition $X_{D_0}(s_0-1)=s_0-1$ in Definition~\ref{def:prepared} is part of the audit's
definition and is not used in the proof, because $\rot$ is applied at every step. The theorem does
not identify $G_{\mathrm{tr}}$ with a grid of the connected sum of the knots of the factors. That
geometric statement is record M613 and remains open.
\end{remark}

\section{Correspondence with the Lean declarations}\label{sec:lean}

All declarations are in the namespace \lean{OnLeConjecture.KnotFloer.GridFull}. The audits
\lean{AuditGridLaurentFiniteTensorAlgebra.lean} and \lean{AuditGridConnectedSumIterated.lean}
check that every declaration of the six modules, including internal ones, depends only on
\lean{propext}, \lean{Classical.choice} and \lean{Quot.sound}.

\begin{longtable}{>{\raggedright\arraybackslash}p{0.34\textwidth}>{\raggedright\arraybackslash}p{0.58\textwidth}}
\toprule
Statement & Lean declaration and module \\
\midrule
Lemma~\ref{lem:pieces} & \lean{PieceComplexData.hom}, \lean{PieceComplexData.iso},
\lean{PieceComplexData.hom_shifts} (\lean{GridLaurentFiniteTensorSingleton}) \\
Proposition~\ref{prop:singleton} & \lean{laurentFiniteMaslovSubsingletonIso} with
\lean{_hom_shifts}, \lean{_inv_shifts}, \lean{_hom_canonicalClass},
\lean{_inv_canonicalClass} (\lean{GridLaurentFiniteTensorSingleton}) \\
Proposition~\ref{prop:reindex} & \lean{laurentFiniteMaslovReindexIso} with \lean{_hom_shifts},
\lean{_inv_shifts}, \lean{_hom_canonicalClass}, \lean{_inv_canonicalClass}, and
\lean{laurentFiniteMaslovFamilyIso} (\lean{GridLaurentFiniteTensorReindex}) \\
Proposition~\ref{prop:rows} & \lean{laurentMaslovCyclicRowsPowIso} with \lean{_hom_shifts},
\lean{_inv_shifts}, \lean{gridCanonicalClass_cyclicRowsPow},
\lean{gridCanonicalClass_cyclicRowsPow_inv} (\lean{GridFullCyclicRowsPower}) \\
Lemma~\ref{lem:insert}, Definitions~\ref{def:laurent-map} and~\ref{def:induced},
Lemma~\ref{lem:bridge} & \lean{insertLinear}, \lean{LaurentGridFiniteMap},
\lean{LaurentTensorFiniteMap}, \lean{LaurentGridInduced}, \lean{LaurentTensorInduced},
\lean{LaurentGridInduced.exists_laurentGridFiniteMap}
(\lean{GridLaurentFiniteTensorReassociation}) \\
Theorem~\ref{thm:reassociation} & \lean{exists_tensorGrid_finiteTensor_map},
\lean{LaurentGridFiniteMap.insertTensor_canonicalClass}
(\lean{GridLaurentFiniteTensorReassociation}) \\
Lemma~\ref{lem:m608} & \lean{seamParity_of_positivePaths}
(\lean{GridConnectedSumSeamUnconditional}) \\
Definition~\ref{def:prepared}, Lemmas~\ref{lem:size} and~\ref{lem:prepare} & \lean{SeamPrepared},
\lean{iteratedSeamSize}, \lean{iteratedSeamGrid}, \lean{two_le_iteratedSeamSize},
\lean{exists_seamPrepared_movesTo} (\lean{GridConnectedSumIterated}) \\
Theorem~\ref{thm:iterated} & \lean{exists_seamChain_finiteTensor_map},
\lean{exists_iteratedSeam_finiteTensor_map_of_parity},
\lean{exists_iteratedSeam_finiteTensor_map} (\lean{GridConnectedSumIterated}) \\
\bottomrule
\end{longtable}

\subsection*{Deviations from the proposed API}
\begin{enumerate}
\item The index types of the isomorphisms of M005, M610 and the family transport are in the
universe \lean{Type}. For $K=\F$ the native grid complex lives in the universe of
\lean{Type}, and a finite tensor over an index type of a higher universe lives in that higher
universe, so the proposed isomorphisms do not typecheck for such index types.
\item M609 carries the hypothesis that $\theta$ is $\Lam$-induced and the same conclusion for
$\psi$, as in the record of M609 in the audit. The proposal cited in the job listing of J03 omits
both (Remark~\ref{rem:linear}). This formulation was approved on 2 November 2026.
\item In addition to the proposed conditional theorem, the module proves the unconditional
version \lean{exists_iteratedSeam_finiteTensor_map}, using M601.
\end{enumerate}

\end{document}
